% 4. Related Work [cite: 43]
\section{Related Work}

% Summarize the most relevant previous work and place your project in context. Organize it by themes, not by just listing papers. [cite: 44, 45]
% For each theme, explain: what the prior work does [cite: 48], what its limitations are [cite: 49], and how your work builds on or differs from it. [cite: 50]

\subsection{Crowd Simulation and Agent-Based Methods}

Traditional agent-based crowd simulations often relied on Social Force Models (SFM), which treat agents as particles influenced by repulsive forces based on physical proximity. However, these models frequently struggle with "jitter" and fail to capture the proactive nature of human navigation. Karamouzas et al. \cite{Karamouzas2014} revolutionized this domain by identifying a universal power law that governs pedestrian interactions based on the time-to-collision ($\tau$) rather than spatial distance.

Their research demonstrates that the interaction energy between pedestrians follows a $1/\tau^2$ relationship, proving that human avoidance behavior is fundamentally anticipatory. Pedestrians do not simply react to those nearby but they proactively adjust their trajectories based on predicted future collisions. This "anticipatory energy" provides a much more robust framework for agent steering, allowing for the natural emergence of self-organized behaviors like lane formation and smooth bottleneck navigation.

This shift from reactive, force-based repulsion to predictive, energy-based optimization serves as a critical precursor to modern Position-Based Dynamics (PBD) approaches. While Karamouzas et al.\cite{Karamouzas2014} defined the statistical mechanics of how agents should move to avoid collisions, PBD provides the computational framework to enforce these anticipatory requirements as hard kinematic constraints, enabling the simulation of high-density crowds with both physical stability and behavioral realism.

\subsection{PBD for Collision Handling and Motion Constraints} % [cite: 46]
While the universal power law provides a mathematical foundation for anticipatory avoidance, implementing these interactions in large-scale simulations requires a solver that is both stable and computationally efficient. Weiss et al. \cite{weiss2019} addressed this by introducing Position-Based Multi-Agent Dynamics (PBMAD). This framework reformulates the time-to-collision power law—originally proposed by Karamouzas et al.\cite{Karamouzas2014}—into a series of position-based constraints within a Position-Based Dynamics (PBD) pipeline.

In PBMAD, the anticipatory interaction energy is treated as a constraint function $C(\mathbf{x})$, which is solved directly at the position level. Unlike force-based methods that require small time steps to avoid numerical instability (the "explosion" of agents), the PBD approach allows for much larger time steps while maintaining strict stability. Weiss et al. \cite{weiss2019} specifically demonstrates that by using a Long-Range Collision Constraint (LRCC), agents can resolve potential collisions well before they occur, mimicking the $1/\tau^2$ behavior in a way that is highly parallelizable. This enables the real-time simulation of tens of thousands of agents, maintaining high-density "flow" patterns without the prohibitive overhead of traditional velocity-based or force-based optimization.

\subsection{Spatial Partitioning and Acceleration Structures} % [cite: 46]
For rigid-body interactions and narrow-phase contact resolution, hierarchical spatial partitioning (e.g., octree combined with AABB pruning) provides an effective defense against quadratic contact-check explosion. A shallow, GPU-traversal-friendly hierarchy—often a two-level scheme—strikes a balance between culling efficiency and traversal overhead. Collision response commonly employs a two-stage strategy: mass-weighted penetration correction (frequently using Baumgarte-style stabilization) to remove interpenetration reliably, followed by impulse-based velocity updates using Coulomb friction models to produce plausible restitution and tangential sliding or sticking behaviors. Such hierarchical frameworks have demonstrated robustness for complex meshes and have enabled interactive rigid-body simulations at scales where flat broad-phase methods fail \cite{sung}.

\subsection{Foveated rendering and Perceptual Optimization} % [cite: 46]
While foveated rendering traditionally focuses on reducing shading and pixel density in the periphery, Stancu et al. \cite{stancu} extend this concept to the underlying simulation pipeline with Foveated Animations. This approach recognizes that the human visual system is less sensitive to high-frequency motion and fine skeletal details in the peripheral field. By utilizing gaze-tracking data, the system dynamically adjusts the animation update rates and the complexity of bone transformations based on an agent's eccentricity from the viewer's focal point. The methodology leverages "visual crowding" a phenomenon where peripheral stimuli are perceived statistically rather than individually.

In large-scale crowd scenarios, the authors demonstrate that the update frequency for character skinning and blending can be significantly reduced in the periphery without a perceptible loss in quality. Their findings indicate that in optimal scenarios, the number of required animation operations can be reduced by up to 99.3\%. This technique is particularly synergistic with Position-Based Dynamics (PBD) solvers. While the PBD solver ensures stable global navigation and collision avoidance, foveated animation handles the local, high-cost task of deforming character meshes. By decoupling the simulation of an agent's position from the frequency of its skeletal updates, developers can achieve "cinematic" crowd densities in VR and head-mounted displays that were previously bottlenecked by the CPU-intensive nature of character skinning.

\subsection{GPU Based Simulation} % [cite: 46]
GPU-based simulation frameworks have become increasingly prevalent for large-scale crowd simulations due to their ability to handle massive parallelism. When it comes to crowd simulation, there are several key areas where GPU acceleration can be particularly beneficial. For instance, the broad-phase collision detection can be efficiently implemented using spatial partitioning techniques such as uniform grids or hierarchical structures like octrees, which can be traversed in parallel on the GPU. In other words, the neigbhor search can be performed in a single pass, significantly reducing the computational overhead compared to CPU-based methods. 

While the previous methods focus primarily on agent navigation, scaling to large crowds of rigid-body entities requires addressing the significant computational overhead of collision detection and response. Sung and Hong \cite{sung} propose a GPU-optimized framework designed to maintain real-time performance for thousands of interacting rigid bodies. Their approach moves beyond simple bounding volumes, utilizing a parallelized pipeline that optimizes both the broad-phase and narrow-phase collision detection stages.

The core of their contribution lies in a scalable collision response mechanism that minimizes the synchronization overhead typically found in GPU-based solvers. By restructuring the contact resolution process to be more cache-friendly and reducing memory contention between threads, the authors demonstrate that rigid-body simulations can scale significantly without the exponential performance decay seen in traditional CPU-based or naive GPU implementations.

In the context of crowd simulation, this rigid-body optimization is essential for scenarios where agents are not merely moving particles but are interacting with complex environments or each other through physical contact—such as in high-density bottleneck scenarios or vehicle-based simulations. When paired with the anticipatory constraints of PBD, the GPU-optimized collision handling by Sung and Hong \cite{sung} ensures that the physical realism of the crowd does not compromise the frame rate.